
% ---------------------------------------------------------------------------
% Case 2: Amazon hiring tool
% ---------------------------------------------------------------------------
\begin{frame}[allowframebreaks]{Case 2: A Resume-Screening Tool (Reuters reporting, 2018)}

    \begin{tcolorbox}[colback=raiAmber!5, colframe=raiAmber!70, title=Note on evidence]
        \small This case comes from \textbf{news reporting based on anonymous sources}, not from a
        published study. Treat the mechanism as illustrative and the details as journalistic.
    \end{tcolorbox}

    \vspace{0.6em}

    According to Reuters, a team at Amazon built experimental tools from 2014 that scored job
    applicants one to five stars. By 2015 the team found the system was not rating candidates for
    software-developer roles in a gender-neutral way.

    \vspace{0.6em}

    \begin{itemize}
        \setlength\itemsep{0.4em}
        \item The models were trained on resumes submitted to the company over roughly a decade ---
              a pool that was predominantly male in technical roles.
        \item The reporting states the system penalised resumes containing the word
              ``women's'' and downgraded graduates of two all-women's colleges.
        \item Amazon said the tool was never used by recruiters to evaluate candidates, and the
              project was abandoned.
    \end{itemize}

    \vspace{0.5em}

    \begin{block}{The mechanism, stated plainly}
        \small The label was not ``a good engineer.'' It was ``a person this company hired
        before.'' Removing the gender field does not help, because the model finds
        \textbf{proxies}. The learning worked; the \textbf{target definition} was the defect.
    \end{block}

    \vspace{0.3em}
    {\scriptsize J.~Dastin, ``Amazon scraps secret AI recruiting tool that showed bias against
    women,'' \emph{Reuters}, 10 October 2018.}
\end{frame}
